\uuid{kizY}
\titre{Espace engendré par une famille de vecteurs/Famille libre}
\auteur{hanani ; mzari}
\organisation{exo7}
\licence{CC BY-NC-SA 4.0}
\datecreate{2020-01-20}
\niveau{L1}
\module{Algèbre}
\chapitre{Espace vectoriel}
\sousChapitre{Base}
\difficulte{2}
\theme{Base et dimension, Espace engendré par une famille de vecteurs, Famille libre}
\exo7qcmid{lille-2.espaces-vectoriels.030}

\contenu{
\begin{question}
\qclassification{106.04}
\qtags{motcle=Espace engendré par une famille de vecteurs/Famille libre}
Dans $\Rr^4$, on considère les vecteurs $u_1=(1,1,-1,0)$, $u_2=(0,1,1,1)$ et $u_3=(1,-1,a,b)$, où $a$ et $b$ sont des réels. Quelles sont les assertions vraies ?
\begin{answers}
\bad{$\forall \, a,b \in \Rr, u_3 \notin \mbox{Vect}\{u_1,u_2\}$}.
\good{$\exists \,  a,b \in \Rr, u_3 \in {\mbox{Vect}} \{u_1,u_2\}$}.
\good{$u_3\in \mbox{Vect}\{u_1,u_2\}$ si et seulement si $a=-3$ et $b=-2$}.
\bad{$\forall \, a,b \in \Rr, \; \{u_1,u_2,u_3\}$ est libre}.
\end{answers}
\begin{explanations} $u_3 \in {\mbox{Vect}} \{u_1,u_2\}$ si, et seulement si, il existe $\alpha, \beta \in \Rr$ tels que
$u_3=\alpha u_1+ \beta u_2$. En résolvant ce système, on obtient $b=-2$ et $a=-3$. Pour $a=-3$ et $b=-2$, la famille $\{u_1,u_2,u_3\}$ n'est pas libre.
\end{explanations}
\end{question}
}
